\chapter{Dense CW Tu Sweep}
\label{ch:cwtu}

\section{Campaign setup}

The dense CW grid was the post-fix measurement intended to produce a clean \(\fres(\Sd)\) curve. Driver: \texttt{scripts/run\_sheath\_cw\_tu.ps1}. Output: \texttt{results/sheath\_cw\_tu/} (voltage/current only; no field dumps).

\begin{itemize}[leftmargin=1.5em]
\item \(\Sd\in\{0,2,4,6,8,10\}\).
\item Drive frequencies: 1.50, 1.60, 1.70, 1.75, 1.80, 1.85, 1.90, 2.00, 2.10, 2.20, 2.30~MHz (66 cases).
\item \(\fp=\SI{2}{\mega\hertz}\), template \texttt{sheath\_sine\_vc.str}, \(\mathrm{MaxParallel}=3\), \(\mathrm{OMP}=4\).
\item Started 13 August 2026 12:48; all 66 cases finished by 18 August 2026.
\end{itemize}

Analyzer: \texttt{scripts/analyze\_cw\_tu\_partial.py} (last-50\% phasor \(Z\), same convention as the \(\fp\)-screen; series resonance is the first \(\ImZ\) \(+\) to \(-\) crossing). Tables and figures: \texttt{paper0_sheath_campaign/analysis/data/cw\_tu\_*.txt}, \texttt{paper0_sheath_campaign/analysis/figures/cw\_tu\_*.png}.

\section{Campaign status}

As of 18 August 2026: \textbf{66 complete, 0 running, 0 pending}, 0 failed. Every planned cell in Table~\ref{tab:cwtu-status} finished with a full \texttt{data.vc} file.

The \(\Sd=0\), 1.50~MHz case was skipped once by the orchestrator on restart; its file is smaller than the others and lacks a \texttt{DONE} line, but the phasor at 1.50~MHz is consistent with the rising inductive branch below the crossing and is included in the plots.

\begin{table}[ht]
\centering
\caption{CW Tu case status (18 August 2026). All cells complete.}
\label{tab:cwtu-status}
\scriptsize
\begin{tabular}{@{}rccccccccccc@{}}
\toprule
\(\Sd\) & 1.50 & 1.60 & 1.70 & 1.75 & 1.80 & 1.85 & 1.90 & 2.00 & 2.10 & 2.20 & 2.30 \\
\midrule
0  & C & C & C & C & C & C & C & C & C & C & C \\
2  & C & C & C & C & C & C & C & C & C & C & C \\
4  & C & C & C & C & C & C & C & C & C & C & C \\
6  & C & C & C & C & C & C & C & C & C & C & C \\
8  & C & C & C & C & C & C & C & C & C & C & C \\
10 & C & C & C & C & C & C & C & C & C & C & C \\
\bottomrule
\end{tabular}
\end{table}

\section{Impedance versus frequency}

Figure~\ref{fig:cwtu-z} and Table~\ref{tab:cwtu-z} show the full-grid phasor results. Unlike the pre-fix \(\fp\)-screen, curves for different \(\Sd\) are well separated across 1.50--2.30~MHz.

\begin{figure}[ht]
\centering
\includegraphics[width=0.92\textwidth]{cw_tu_impedance.png}
\caption{CW Tu sweep, all 66 cases. \(\Sd=0\) crosses \(\ImZ=0\) near 1.85~MHz and is capacitive above \(\sim\SI{2.0}{\mega\hertz}\). Every \(\Sd\ge 2\) curve stays capacitive (\(\ImZ<0\)) across the entire band.}
\label{fig:cwtu-z}
\end{figure}

\begin{table}[ht]
\centering
\caption{Selected phasor impedances (\(\si{\ohm}\)). Full table: \texttt{paper0_sheath_campaign/analysis/data/cw\_tu\_summary.txt}.}
\label{tab:cwtu-z}
\begin{tabular}{@{}rrrrr@{}}
\toprule
\(f\) (MHz) & \(\Sd\) & \(\ReZ\) & \(\ImZ\) & note \\
\midrule
1.50 & 0  & 2130  & $+585$   & inductive shoulder \\
1.85 & 0  & 62775 & $-640$   & just past crossing \\
2.00 & 0  & 54406 & $-54920$ & capacitive above \(\fres\) \\
1.50 & 2  & 15227 & $-9181$  & capacitive at band edge \\
2.00 & 2  & 21484 & $-40907$ & still capacitive \\
1.70 & 10 & 6182  & $-20959$ & matches confirm-Z \\
2.00 & 10 & 2238  & $-22032$ & \(\lvert Z\rvert\) lowest in band \\
\bottomrule
\end{tabular}
\end{table}

\section{\texorpdfstring{Series resonance $\fres(\Sd)$}{Series resonance f\_res(Sd)}}

Linear interpolation of the \(\Sd=0\) \(+\) to \(-\) crossing between 1.80 and 1.85~MHz gives
\begin{equation}
\fres(\Sd=0)\approx\SI{1.846}{\mega\hertz}.
\label{eq:cwtu-fres0}
\end{equation}
That matches the post-fix confirm (sign change between 1.7 and 1.9~MHz) and is slightly above the coarse pre-fix 1.7/1.9 interpolation (\(\sim\SI{1.775}{\mega\hertz}\)).

For \(\Sd=2,4,6,8,10\), \(\ImZ<0\) at \emph{every} frequency in 1.50--2.30~MHz. No \(+\) to \(-\) crossing appears anywhere in the grid for \(\Sd>0\). Any series crossing for those cases therefore lies \emph{below} \(\SI{1.50}{\mega\hertz}\)---the opposite direction from Tu~\cite{tu2008}, which predicts \(\fres\) moving \emph{upward} with \(\Sd\) toward the free-space value. Figure~\ref{fig:cwtu-res} contains only the \(\Sd=0\) point.

\begin{figure}[ht]
\centering
\includegraphics[width=0.85\textwidth]{cw_tu_resonance.png}
\caption{\(\fres(\Sd)\) from the full CW grid (\(\ImZ\) \(+\) to \(-\) crossing). Only \(\Sd=0\) has a crossing in 1.50--2.30~MHz.}
\label{fig:cwtu-res}
\end{figure}

Table~\ref{tab:cwtu-res} summarizes the resonance extraction. The high-frequency extension (1.90--2.30~MHz) rules out the provisional hypothesis from the partial run that a Tu-like upward crossing might appear above 1.9~MHz for large \(\Sd\): \(\Sd=10\) remains strongly capacitive at 2.0 and 2.3~MHz.

\begin{table}[ht]
\centering
\caption{Series resonance from \(\ImZ\) \(+\) to \(-\) crossings.}
\label{tab:cwtu-res}
\begin{tabular}{@{}rll@{}}
\toprule
\(\Sd\) & \(\fres\) & note \\
\midrule
0  & \(\SI{1.846}{\mega\hertz}\) & crossing in 1.50--2.30~MHz band \\
2  & --- & all \(\ImZ<0\); crossing at or below 1.50~MHz \\
4  & --- & all \(\ImZ<0\); crossing at or below 1.50~MHz \\
6  & --- & all \(\ImZ<0\); crossing at or below 1.50~MHz \\
8  & --- & all \(\ImZ<0\); crossing at or below 1.50~MHz \\
10 & --- & all \(\ImZ<0\); crossing at or below 1.50~MHz \\
\bottomrule
\end{tabular}
\end{table}

\section{What succeeded}

\begin{itemize}[leftmargin=1.5em]
\item Coupling is visible in CW at every \(\Sd>0\): \(Z(f)\) is no longer independent of sheath width.
\item The \(\Sd=0\) plasma-loaded resonance is resolved to \(\sim\SI{50}{\kilo\hertz}\) spacing near \(\fp\).
\item The full grid confirms the post-fix confirm-Z picture: \(\Sd=0\) is inductive below \(\sim\SI{1.85}{\mega\hertz}\); \(\Sd=10\) is capacitive everywhere in the band.
\end{itemize}

\section{What did not match the upward-shift target}

The campaign design target---inspired by a circuit reading of Tu~\cite{tu2008}---was an \emph{upward} shift of \(\fres\) with \(\Sd\). That shift was not reproduced in this band:

\begin{itemize}[leftmargin=1.5em]
\item Only \(\Sd=0\) yields a measurable in-band series crossing.
\item Wider sheaths add capacitive loading (\(\lvert Z\rvert\) falls with \(\Sd\) at fixed \(f\)) but do not move a \(+\) to \(-\) crossing upward through 1.50--2.30~MHz.
\item If the same crossing definition applies to \(\Sd\ge 2\), the feature has moved \emph{downward} relative to \(\Sd=0\), or the feed impedance is dominated by series sheath capacitance without a clear in-band LC crossing.
\end{itemize}

Follow-up completed in Chapter~\ref{ch:lowf}: the CW grid was extended below 1.50~MHz for \(\Sd=0,2,10\). That campaign locates \(\fres(\Sd=2)\approx\SI{0.655}{\mega\hertz}\) and finds \(\Sd=10\) capacitive down to \(\SI{0.50}{\mega\hertz}\), confirming the downward trend. Interpretation (series \(C_{\mathrm{sh}}\) vs dielectric unloading; Balmain--Liu lineage; Tu regime mismatch): \texttt{paper0_sheath_campaign/analysis/tu\_discrepancy\_discussion.md} and Chapter~\ref{ch:status}.
